
本章将证明第~7 章中讨论的\BookRecordIndex{T0143}格罗莫夫--威滕不变量\BookRecordIndex{T0125}分裂公理（定理~\ref{thm:7-5-10}）。
其基础是如下观察：两条以适当方式横截相交的 $J$-全纯球面可以\BookRecordIndex{T0254}粘合起来，从而产生一族
代表原球面同调类之和的 $J$-全纯球面。这在某种意义上是\BookRecordIndex{T0139}格罗莫夫紧性的逆过程。为了证明
它，首先用\BookRecordIndex{T0227}连通和构造一条近似 $J$-全纯球面，再用\BookRecordIndex{T0380}隐函数定理确定附近确有一条真正的
$J$-全纯球面。在\BookRecordIndex{T0334}伪全纯曲线的背景下，最早的此类\BookRecordIndex{T0256}粘合定理由 Floer
\srcite{113,114,116} 给出；他考察的是连接两条周期轨道的轨线之\BookRecordIndex{T0246}模空间。由于假定周期轨道
非退化，这一情形略为容易。$J$-全纯球面的\BookRecordIndex{T0256}粘合定理最早由 Ruan--Tian \srcite{345}
证明，随后 Liu \srcite{248} 又用与我们略有不同的方法给出证明。这里的论证是本书第一版
论证的改进。后来又出现了许多其他版本，例如见 Fukaya--Ono \srcite{127}、
Ionel--Parker \srcite{201} 和 Zinger \srcite{426}。

第一节建立记号并陈述主要\BookRecordIndex{T0256}粘合定理（定理~\ref{thm:10-1-2}）。随后的六节用于证明它。
第~\ref{sec:10-8} 节证明分裂定理~\ref{thm:7-5-10}。最后，在定理~\ref{thm:10-9-1} 中，我们用
主要\BookRecordIndex{T0256}粘合定理描述紧化\BookRecordIndex{T0246}模空间 $\overline{\mathcal M}_{0,4}(A,J)$ 中单\BookRecordIndex{T0191}节点单纯\BookRecordIndex{T0340}稳定映射之集
的一个邻域结构。我们还会说明，这个简化的\BookRecordIndex{T0256}粘合定理足以建立量子乘法的结合律。最后两节
只用到\BookRecordIndex{T0256}粘合定理的陈述而不用其证明，因此读完第~\ref{sec:10-1} 节便可立即阅读它们。

定理~\ref{thm:10-1-2} 之证明的基本想法，是先用\BookRecordIndex{T0188}截断函数拼合与所求对象非常接近的对象，
再用 Newton 型迭代得到精确解。第一步是从两条相交的 $J$-全纯球面构造一族近似
$J$-全纯球面 $u^R$（$R$ 很大）。这一步称为\emph{\BookRecordIndex{T0394}预粘合}，将在
第~\ref{sec:10-2} 节讨论。下一节为沿\BookRecordIndex{T0395}预粘合曲线 $u^R$ 的向量场和 $1$-形式引入合适的
\BookRecordIndex{T0179}加权索伯列夫范数。为应用第~\ref{sec:3-5} 节的\BookRecordIndex{T0380}隐函数定理，我们需要\BookRecordIndex{T0346}线性化算子
$D_{u^R}$ 的适当\BookRecordIndex{T0392}右逆关于这些\BookRecordIndex{T0177}加权范数的一致估计。依循 Donaldson--Kronheimer
\srcite{88}，我们先构造\BookRecordIndex{T0196}近似右逆 $T_{u^R}$，再得到与 $T_{u^R}$ 有相同像的\BookRecordIndex{T0197}精确右逆
$Q_{u^R}$。

为使这一方案成立，必须使用性质精细的\BookRecordIndex{T0188}截断函数 $\beta$；这些性质在
第~\ref{sec:10-4} 节建立。主要构造见第~\ref{sec:10-5} 节。在那里，我们定义近似逆族
$T_{u^R}$ 和真正的逆族 $Q_{u^R}$，再应用逆函数定理定义\BookRecordIndex{T0257}粘合映射
$\iota_c^{\delta,R}$。我们还会建立它的大部分性质，把一些细节留到技术性的
第~\ref{sec:10-6} 节。定理~\ref{thm:10-1-2} 证明的剩余部分涉及满射性。要证明这一点，
必须理解\BookRecordIndex{T0254}粘合在哪种意义下是\BookRecordIndex{T0192}紧性的逆过程；为此将用到引理~\ref{lem:4-7-3} 关于低能量长
圆柱性质的结论。细节见第~\ref{sec:10-7} 节。第~\ref{sec:10-6}、\ref{sec:10-7} 节
不可避免地较为技术化。为了把握构造的全貌，读者不妨先集中阅读
第~\ref{sec:10-1}、\ref{sec:10-2} 和~\ref{sec:10-5} 节。

\section{\BookRecordIndex{T0256}粘合定理}\label{sec:10-1}

\BookRecordIndex{T0003}$J$-全纯曲线的\BookRecordIndex{T0256}粘合定理可以看作\BookRecordIndex{T0139}格罗莫夫紧性定理的逆。粗略地说，它断言：若两条（或
更多条）\BookRecordIndex{T0003}$J$-全纯曲线相交并满足适当的\BookRecordIndex{T0166}横截性条件，那么代表其同调类之和的一列
\BookRecordIndex{T0003}$J$-全纯曲线可以在\BookRecordIndex{T0141}格罗莫夫收敛的意义下（见定义~\ref{def:5-2-1}）逼近它们。我们的证明
改造了 Donaldson--Kronheimer \srcite{88, pp.~287--295} 为四维流形上的反自对偶瞬子建立\BookRecordIndex{T0256}粘合
定理的方法。要精确陈述定理还需作一些准备。我们将在第~\ref{sec:6-7} 节的框架下描述
\BookRecordIndex{T0257}粘合映射，其中 $J$ 依赖于 $z\in S^2$。因为 $J$ 在一个紧族中变化，这对分析并无本质
影响，却使结果适用得更广。当 $J$ 与 $z$ 无关时会有额外的对称性，从而得到一个略为
简单、可用\BookRecordIndex{T0340}稳定映射和\BookRecordIndex{T0141}格罗莫夫收敛语言陈述的\BookRecordIndex{T0256}粘合定理版本。第~\ref{sec:10-9} 节将讨论
这个版本，并说明如何用它证明最简单版本的分裂原理。

\paragraph{\BookRecordIndex{T0226}连通对的\BookRecordIndex{T0246}模空间。}
如上所述，设
$J^0=\{J_z^0\}_{z\in S^2}$ 和 $J^\infty=\{J_z^\infty\}_{z\in S^2}$
是 $M$ 上两族光滑的 $\omega$-\BookRecordIndex{T0372}驯服\BookRecordIndex{T0088}殆复结构，并设常数 $\kappa>0$ 满足
\begin{equation}\label{eq:10-1-1}
  |z|<\kappa
  \quad\Longrightarrow\quad
  J_z^0=J_{1/\bar z}^{\infty}=:J.
\end{equation}
固定两个同调类 $A^0,A^\infty\in H_2(M;\mathbb Z)$。回忆，
$\mathcal M(A^0;\{J_z^0\})$ 表示 $M$ 中代表类 $A^0$ 的
$\{J_z^0\}$-全纯曲线空间；$\mathcal M(A^\infty;\{J_z^\infty\})$ 的含义类似。
愿意的话，读者可以只考虑与 $z$ 无关的 $J$，但此时应只考察单纯映射。

\begin{definition}\label{def:10-1-1}
称
$(J^0,J^\infty)\in\mathcal J_\tau(S^2;M,\omega)\times
\mathcal J_\tau(S^2;M,\omega)$ 对 $(A^0,A^\infty)$ \emph{正则}，如果
$J^0\in\mathcal J_{\mathrm{reg}}(S^2;A^0)$ 和
$J^\infty\in\mathcal J_{\mathrm{reg}}(S^2;A^\infty)$ 对某个 $\kappa>0$
满足~\eqref{eq:10-1-1}，并且\BookRecordIndex{T0283}求值映射
\begin{equation}\label{eq:10-1-2}
  \begin{aligned}
  \operatorname{ev}:\mathcal M(A^0;J^0)\times\mathcal M(A^\infty;J^\infty)
  &\longrightarrow M\times M,\\
  \operatorname{ev}(u^0,u^\infty)&:=(u^0(0),u^\infty(\infty))
  \end{aligned}
\end{equation}
横截于对角线。所有这样的\BookRecordIndex{T0406}正则对之集记作
$\mathcal J_{\mathrm{reg}}(S^2;A^{0,\infty})$。
\end{definition}

由命题~\ref{prop:6-7-7}，集合
$\mathcal J_{\mathrm{reg}}(S^2;A^0)$ 和
$\mathcal J_{\mathrm{reg}}(S^2;A^\infty)$ 在所有满足~\eqref{eq:10-1-1} 的
$\omega$-\BookRecordIndex{T0372}驯服族 $\{J_z\}\in\mathcal J_\tau(S^2;M,\omega)$ 的空间中都是\BookRecordIndex{T0300}剩余集。
此外，定理~\ref{thm:6-3-1} 的证明表明，对于通有的偶对 $(J^0,J^\infty)$，\BookRecordIndex{T0283}求值映射
\eqref{eq:10-1-2} 横截于对角线。因此，
$\mathcal J_{\mathrm{reg}}(S^2;A^{0,\infty})$ 在所有满足~\eqref{eq:10-1-1} 的对
\[
  (J^0,J^\infty)\in\mathcal J_\tau(S^2;M,\omega)\times
  \mathcal J_\tau(S^2;M,\omega)
\]
中是\BookRecordIndex{T0300}剩余集。再者，若
$(J^0,J^\infty)\in\mathcal J_{\mathrm{reg}}(S^2;A^{0,\infty})$，则\BookRecordIndex{T0246}模空间
\[
  \mathcal M(A^{0,\infty};J^{0,\infty})
  :=\bigl\{(u^0,u^\infty)\in
  \mathcal M(A^0;J^0)\times\mathcal M(A^\infty;J^\infty)
  \bigm|u^0(0)=u^\infty(\infty)\bigr\}
\]
是维数为 $2n+2c_1(A^0+A^\infty)$ 的光滑有限维流形。

给定常数 $c>0$，用
\[
  \mathcal M(c):=\mathcal M(A^{0,\infty};J^{0,\infty},c)
\]
表示所有满足
\[
  \|du^0\|_{L^\infty}\le c,
  \qquad
  \|du^\infty\|_{L^\infty}\le c
\]
的\BookRecordIndex{T0226}连通对 $(u^0,u^\infty)\in\mathcal M(A^{0,\infty};J^{0,\infty})$ 所成的紧子集。
在定理~\ref{thm:10-1-2} 中，我们将考察第一球面上的\BookRecordIndex{T0283}求值映射
\begin{equation}\label{eq:10-1-3}
  \begin{aligned}
  \operatorname{ev}^0:\mathcal M(c)\times
  \bigl(S^2\setminus\operatorname{int}B_{1/c}(0)\bigr)&\longrightarrow M,\\
  \operatorname{ev}^0(u^0,u^\infty,z)&:=u^0(z).
  \end{aligned}
\end{equation}
以下讨论中，$c$ 始终视为任意大的固定常数。

\paragraph{\BookRecordIndex{T0256}粘合定理的陈述。}
以下固定同调类 $A^0,A^\infty\in H_2(M;\mathbb Z)$，并记
\[
  A:=A^0+A^\infty.
\]
还固定\BookRecordIndex{T0406}正则对
$(J^0,J^\infty)\in\mathcal J_{\mathrm{reg}}(S^2;A^{0,\infty})$。
我们将证明：对每个 $c>0$ 和每个充分大的 $R>0$，\BookRecordIndex{T0246}模空间
$\mathcal M(A^{0,\infty};J^{0,\infty},c)$ 都可嵌入\BookRecordIndex{T0246}模空间
$\mathcal M(A;J^R)$，其中
\begin{equation}\label{eq:10-1-4}
  J_z^R:=
  \begin{cases}
    J_z^0,& |z|\ge 1/R,\\
    J_{R^2z}^\infty,& |z|\le 1/R.
  \end{cases}
\end{equation}
这个 $J^R$ 的公式展示了\BookRecordIndex{T0254}粘合如何作用于定义域：
$\mathcal M(A;J^R)$ 的每个元素都是从黎曼球面
$S^2=\mathbb C\cup\{\infty\}$ 到 $M$ 的映射，其中圆盘 $\{|z|\ge1/R\}$ 应与
$u^0$ 定义域中的相应圆盘等同，而圆盘 $\{|z|\le1/R\}$ 在作重标度
$z\mapsto R^2z$ 后，应与 $u^\infty$ 定义域中的圆盘 $\{|z|\le R\}$ 等同。
因此，$R$ 是一个重标度参数（见注~\ref{rem:10-2-4}）。

\BookRecordIndex{T0257}粘合映射 $\mathcal M(c)\to\mathcal M(A;J^R)$ 还依赖另一参数 $\delta$，其倒数对应于
“\BookRecordIndex{T0147}颈部的长度”。式~\eqref{eq:10-5-4} 最清楚地展示了它的作用。三个常数 $c,R,\delta$ 的关系如下：
给定大常数 $c>0$，先固定小常数 $\delta<\delta_0(c)$，再考察所有充分大的 $R$。事实上，
对 $R$ 的有关条件形如 $\delta R>1/\delta_0(c)$；$\delta$ 与 $R$ 确定了\BookRecordIndex{T0170}环域
$\delta/R\le|z|\le1/(\delta R)$，重要的分析将在其上进行。对每个 $\delta_0>0$，记
\BookRecordIndex{T0171}环域参数集为
\begin{equation}\label{eq:10-1-5}
  \mathcal A(\delta_0)
  :=\bigl\{(\delta,R)\bigm|0<\delta<\delta_0,\ \delta R>1/\delta_0\bigr\}.
\end{equation}
下一个定理陈述所得\BookRecordIndex{T0257}粘合映射
\[
  \iota_c^{\delta,R}:\mathcal M(c)\longrightarrow\mathcal M(A;J^R)
\]
的主要性质。由 (iv)，对每个 $(\delta,R)\in\mathcal A(\delta_0)$，它都是保定向的光滑
嵌入。性质 (i) 说明\BookRecordIndex{T0257}粘合映射光滑地依赖于 $R$；性质 (ii) 说明\BookRecordIndex{T0254}粘合曲线
$\widetilde u^R$ 在 $R\to\infty$ 时收敛到它的两个组成部分 $(u^0,u^\infty)$；(v) 则
说明，它的像包含所有经适当重标度后充分接近 $\mathcal M(c-1)$ 的 $J^R$-全纯映射。
最后，(iii) 说明\BookRecordIndex{T0283}求值映射在 $C^1$ 拓扑中收敛。（注意，$C^0$ 收敛由 (ii) 得出；但是要
控制 $\operatorname{ev}^R$ 与 $M$ 中给定循环的\BookRecordIndex{T0166}横截性，还需要 $C^1$ 收敛；见
第~\ref{sec:10-8} 节。）

\begin{theorem}[\BookRecordIndex{T0254}粘合]\label{thm:10-1-2}
固定两个同调类 $A^0,A^\infty\in H_2(M;\mathbb Z)$、一个\BookRecordIndex{T0088}殆复结构\BookRecordIndex{T0406}正则对
$(J^0,J^\infty)\in\mathcal J_{\mathrm{reg}}(S^2;A^{0,\infty})$ 以及常数 $c>1$。
如上记 $\mathcal M(c):=\mathcal M(A^{0,\infty};J^{0,\infty},c)$。则存在常数
$\delta_0>0$、$\varepsilon>0$，并且对每个
$(\delta,R)\in\mathcal A(\delta_0)$ 有映射
\[
  \iota_c^{\delta,R}=\iota_c^R:
  \mathcal M(c)\longrightarrow\mathcal M(A;J^R),
\]
使得下列条件成立。

\begin{enumerate}[label=\textnormal{(\roman*)}]
\item 固定常数 $0<\delta<\delta_0$。则映射
\[
  (1/(\delta\delta_0),\infty)\times\mathcal M(c)
  \longrightarrow C^\infty(S^2,M),
  \qquad (R,u^0,u^\infty)\longmapsto\iota_c^R(u^0,u^\infty)
\]
光滑，并且对每个 $R>1/(\delta\delta_0)$，映射
$\iota_c^R:\mathcal M(c)\to\mathcal M(A;J^R)$ 都是嵌入。

\item 固定常数 $0<\delta<\delta_0$。设
$(u_\nu^0,u_\nu^\infty)$ 是 $\mathcal M(c)$ 中在 $C^\infty$ 拓扑下收敛到
$(u^0,u^\infty)\in\mathcal M(c)$ 的序列，并设
$R_\nu>1/(\delta\delta_0)$ 趋于无穷。记
\[
  \widetilde u_\nu:=\iota_c^{\delta,R_\nu}(u_\nu^0,u_\nu^\infty).
\]
则当 $\nu\to\infty$ 时，$\widetilde u_\nu(z)$ 在
$S^2\setminus\{0\}$ 的紧子集上一致连同各阶导数收敛到 $u^0(z)$，而
$\widetilde u_\nu(z/R_\nu^2)$ 在 $S^2\setminus\{\infty\}$ 的紧子集上一致连同各阶
导数收敛到 $u^\infty(z)$。

\item 固定常数 $0<\delta<\delta_0$；对 $R>1/(\delta\delta_0)$，令
\[
  \operatorname{ev}^R:\mathcal M(A;J^R)\times S^2\longrightarrow M
\]
为自然\BookRecordIndex{T0283}求值映射。则复合
\[
  \operatorname{ev}^R\circ(\iota_c^{\delta,R}\times\operatorname{id}):
  \mathcal M(c)\times\bigl(S^2\setminus\operatorname{int}B_{1/c}(0)\bigr)
  \longrightarrow M
\]
当 $R\to\infty$ 时在 $C^1$ 拓扑中收敛到\BookRecordIndex{T0283}求值映射~\eqref{eq:10-1-3}。

\item 对每个 $(\delta,R)\in\mathcal A(\delta_0)$，嵌入
$\iota_c^{\delta,R}$ 都保定向。

\item 若 $(\delta,R)\in\mathcal A(\delta_0)$、
$(u^0,u^\infty)\in\mathcal M(c-1)$，并且
$v\in\mathcal M(A;J^R)$ 满足
\[
  \sup_{|z|\ge1}d(v(z),u^0(z))<\varepsilon,
  \qquad
  \sup_{|z|\le1}d(v(z/R^2),u^\infty(z))<\varepsilon,
\]
则 $v$ 属于 $\iota_c^{\delta,R}$ 的像。
\end{enumerate}
\end{theorem}

\begin{remark}\label{rem:10-1-3}
\begin{enumerate}[label=\textnormal{(\roman*)}]
\item 定理~\ref{thm:10-1-2} 证明中的构造表明，映射
$\iota_c^{\delta,R}:\mathcal M(c)\to\mathcal M(A;J^R)$ 与 $c$ 无关。换言之，若
$0<c<c'$，并把 $\delta_0>0$ 取得充分小，则对每个
$(\delta,R)\in\mathcal A(\delta_0)$，映射 $\iota_c^{\delta,R}$ 与
$\iota_{c'}^{\delta,R}$ 在共同定义域 $\mathcal M(c)$ 上一致。

\item 我们的构造在适当重标度的意义下关于两个分支 $u^0,u^\infty$ 对称。因此，
定理~\ref{thm:10-1-2}(iii) 蕴含如下结论。考察由
\[
  \operatorname{ev}^\infty(u^0,u^\infty,z):=u^\infty(z)
\]
定义的\BookRecordIndex{T0283}求值映射
$\operatorname{ev}^\infty:\mathcal M(c)\times B_c(0)\to M$。若记
$u_c^{\delta,R}:=\iota_c^{\delta,R}(u^0,u^\infty)$，则\BookRecordIndex{T0283}求值映射
\[
  \mathcal M(c)\times B_c(0)\longrightarrow M,
  \qquad (u^0,u^\infty,z)\longmapsto u_c^{\delta,R}(z/R^2)
\]
在 $R\to\infty$ 时于 $C^1$ 拓扑中收敛到 $\operatorname{ev}^\infty$。
\end{enumerate}
\end{remark}

定理~\ref{thm:10-1-2}(v) 的满射性陈述略显微妙。对依赖于
$z\in S^2$ 的通有 $J^\infty$，一个映射 $v$ 至多对一个 $R$ 是 $J^R$-全纯的。然而，
若这一条件不成立，则每个 $v$ 都会对某个 $R$ 值的开子集满足 (v) 的条件，而 (v) 蕴含
\BookRecordIndex{T0257}粘合映射 $\iota_c^{\delta,R}$ 的像会随 $R$ 改变而彼此重叠。还要注意，对
$v\in\operatorname{im}\iota_c^{\delta,R}$，$|dv(0)|$ 被常数乘 $R^2$ 控制。因此，对固定
$R$，$\iota_c^{\delta,R}$ 的像不包含一整列\BookRecordIndex{T0141}格罗莫夫收敛序列；如下推论所示，为捕捉
这样的整个序列，必须允许 $R\to\infty$。第~\ref{sec:10-9} 节将对与 $z$ 无关的\BookRecordIndex{T0088}殆复
结构给出一个涉及所有 $R$ 值的\BookRecordIndex{T0256}粘合定理版本，并把 $R$ 的选取与四个\BookRecordIndex{T0066}标记点的\BookRecordIndex{T0185}交比联系起来。

\begin{corollary}\label{cor:10-1-4}
设 $A^0,A^\infty\in H_2(M;\mathbb Z)$，且
$(J^0,J^\infty)\in\mathcal J_{\mathrm{reg}}(S^2;A^{0,\infty})$。固定两个常数
$c>1$ 与 $0<\delta<\delta_0(c)$。设 $R_\nu\to\infty$，且
$u_\nu\in\mathcal M(A;J^{R_\nu})$ 是在 $S^2\setminus\{0\}$ 的紧子集上一致连同各阶
导数收敛到 $u^0(z)$ 的序列；再假设其重标度
$z\mapsto u_\nu(z/R_\nu^2)$ 在 $S^2\setminus\{\infty\}$ 的紧子集上一致连同各阶
导数收敛到 $u^\infty(z)$，其中 $(u^0,u^\infty)\in\mathcal M(c-1)$。则存在常数
$\nu_0$，使得对 $\nu\ge\nu_0$，都有
$u_\nu\in\operatorname{im}\iota_c^{\delta,R_\nu}$。
\end{corollary}

\begin{proof}
当 $\nu$ 充分大时，二元组 $(u_\nu,R_\nu)$ 满足定理~\ref{thm:10-1-2}(v) 的假设。
\end{proof}

定理~\ref{thm:10-1-2} 的证明策略如下。给定两条曲线
$u^0\in\mathcal M(A^0;J^0)$ 和 $u^\infty\in\mathcal M(A^\infty;J^\infty)$，且
$u^0(0)=u^\infty(\infty)$，我们将考察一族近似 $J^R$-全纯曲线
$u^R=u^0\mathbin{\#_R}u^\infty$（$R$ 充分大），并证明其附近存在真正的
$J^R$-全纯曲线 $\widetilde u^R$。论证中需要证明\BookRecordIndex{T0128}弗雷德霍姆算子 $D_{u^R}$ 满射，并估计
它的一个\BookRecordIndex{T0392}右逆；这些步骤依赖于关于\BookRecordIndex{T0188}截断函数的一些初等而精细的观察。随后，要验证所得映射
$\iota_c^R$ 具有上述 (i)--(iii) 性质相当容易。论证的这些步骤在
第~\ref{sec:10-5} 节完成。建立 (iv)、(v) 还需若干进一步估计，但没有本质上的新思想。
为了便于读者理解要点，我们把这部分论证留到第~\ref{sec:10-6}、\ref{sec:10-7} 节。
